\chapter{Measures}

\section{Sigma-algebras and generated structures}

\begin{de}[Algebras and sigma-algebras]
Let $X$ be a set. An \emph{algebra} $\cA\subseteq\pow(X)$ contains $\varnothing$ and is closed under complements and finite unions. A \emph{sigma-algebra} $\cM\subseteq\pow(X)$ is an algebra that is closed under countable unions. A pair $(X,\cM)$ is called a \emph{measurable space}.
\end{de}

\begin{remark}
Because an algebra is closed under complements, it is a sigma-algebra if and only if it is closed under countable intersections. This is an immediate consequence of De Morgan's laws.
\end{remark}

\begin{eg}
The power set $\pow(X)$ and the trivial sigma-algebra $\{\varnothing,X\}$ are sigma-algebras. If $X$ is uncountable, the collection of all countable or co-countable subsets of $X$ is also a sigma-algebra.
\end{eg}

\begin{de}[Generated sigma-algebra]
For $\cE\subseteq\pow(X)$, the sigma-algebra generated by $\cE$ is
\[
\sigma(\cE)=\bigcap\{\cM:\cM\text{ is a sigma-algebra on }X\text{ and }\cE\subseteq\cM\}.
\]
\end{de}

\begin{lem}[Generator principle]
Let $\cE\subseteq\pow(X)$ and let $\cM$ be a sigma-algebra on $X$. If $\cE\subseteq\cM$, then $\sigma(\cE)\subseteq\cM$.
\end{lem}

\begin{de}[Borel sigma-algebra]
If $X$ is a topological space, the \emph{Borel sigma-algebra} $\cB_X$ is the sigma-algebra generated by the open subsets of $X$.
\end{de}

\begin{prop}[Generators of $\cB_{\R}$]
The Borel sigma-algebra on $\R$ is generated by any one of the following families:
\begin{enumerate}
  \item open intervals $(a,b)$;
  \item closed intervals $[a,b]$;
  \item half-open intervals $(a,b]$ or $[a,b)$;
  \item open rays $(a,\infty)$ or $(-\infty,a)$;
  \item closed rays $[a,\infty)$ or $(-\infty,a]$.
\end{enumerate}
It is enough in each case to restrict the endpoints to rational numbers whenever the resulting family still separates all real endpoints by countable unions or intersections.
\end{prop}

\begin{proof}
Every displayed set is Borel, so the sigma-algebra generated by any listed family is contained in $\cB_{\R}$. Conversely, open intervals can be obtained from any of the other families by countable unions, intersections, and complements. Every open subset of $\R$ is a countable union of open intervals with rational endpoints, so each listed family generates all Borel sets.
\end{proof}

\begin{de}[Product sigma-algebra]
Let $(X_\alpha,\cM_\alpha)_{\alpha\in A}$ be measurable spaces and let $X=\prod_{\alpha\in A}X_\alpha$. The product sigma-algebra
\[
\bigotimes_{\alpha\in A}\cM_\alpha
\]
is generated by the coordinate cylinders $\pi_\alpha^{-1}(E)$, where $E\in\cM_\alpha$ and $\pi_\alpha:X\to X_\alpha$ is the coordinate projection.
\end{de}

\begin{prop}
If $A$ is countable, then $\bigotimes_{\alpha\in A}\cM_\alpha$ is generated by all measurable rectangles $\prod_{\alpha\in A}E_\alpha$ with $E_\alpha\in\cM_\alpha$.
\end{prop}

\begin{proof}
Every cylinder is a rectangle in which all but one coordinate set is the whole factor. Conversely, when $A$ is countable,
\[
\prod_{\alpha\in A}E_\alpha
 =\bigcap_{\alpha\in A}\pi_\alpha^{-1}(E_\alpha),
\]
a countable intersection of cylinders.
\end{proof}

\begin{prop}[Generators in the factors]
Suppose $\cM_\alpha=\sigma(\cE_\alpha)$ for every $\alpha$. Then $\bigotimes_{\alpha\in A}\cM_\alpha$ is generated by the cylinders $\pi_\alpha^{-1}(E)$ with $E\in\cE_\alpha$.
\end{prop}

\begin{proof}
Let $\cF$ be the sigma-algebra generated by those cylinders. For fixed $\alpha$, define
\[
\cC_\alpha=\{B\subseteq X_\alpha:\pi_\alpha^{-1}(B)\in\cF\}.
\]
Then $\cC_\alpha$ is a sigma-algebra containing $\cE_\alpha$, hence it contains $\cM_\alpha$. Thus every measurable cylinder belongs to $\cF$, which proves the claim.
\end{proof}

\begin{prop}[Borel sigma-algebras of finite products]
Let $X_1,\dots,X_n$ be metric spaces and let $X=\prod_{j=1}^nX_j$ have the product topology. Then
\[
\bigotimes_{j=1}^n\cB_{X_j}\subseteq\cB_X.
\]
If all $X_j$ are separable, equality holds.
\end{prop}

\begin{proof}
Each coordinate projection is continuous, so every Borel cylinder is Borel in $X$. This proves the first inclusion.

If each $X_j$ is separable, choose a countable base $\cU_j$ for $X_j$. The finite products $U_1\times\cdots\times U_n$, with $U_j\in\cU_j$, form a countable base for $X$. Every open set is therefore a countable union of measurable rectangles and belongs to the product sigma-algebra.
\end{proof}

\begin{cor}
For every $n\geq1$,
\[
\cB_{\R^n}=\bigotimes_{j=1}^n\cB_{\R}.
\]
\end{cor}

\begin{de}[Elementary family]
A collection $\cE\subseteq\pow(X)$ is an \emph{elementary family} if
\begin{enumerate}
  \item $\varnothing\in\cE$;
  \item $E,F\in\cE$ implies $E\cap F\in\cE$;
  \item for every $E\in\cE$, the complement $E^c$ is a finite disjoint union of members of $\cE$.
\end{enumerate}
\end{de}

\begin{prop}
The collection $\cA$ of finite disjoint unions of members of an elementary family $\cE$ is an algebra.
\end{prop}

\begin{proof}
Finite intersections of members of $\cA$ are obtained by distributing intersections over their finite decompositions, and each resulting piece belongs to $\cE$. Complements are obtained by expanding the elementary complement decompositions and refining the finitely many pieces into a disjoint family. Hence $\cA$ is closed under complements and finite intersections, and therefore under finite unions.
\end{proof}

\begin{prop}[Infinite sigma-algebras]
Every infinite sigma-algebra contains a sequence of pairwise disjoint nonempty sets. Consequently, every infinite sigma-algebra has cardinality at least $2^{\aleph_0}$.
\end{prop}

\begin{proof}
Call a nonempty measurable set an \emph{atom} if it contains no nonempty proper measurable subset. If there are infinitely many atoms, choose a countable subfamily. If there are only finitely many atoms, their union cannot be all of $X$, for otherwise every measurable set would be a union of those finitely many atoms and the sigma-algebra would be finite. The complement therefore contains a nonempty atomless measurable set. Repeatedly split the remaining atomless set into two nonempty measurable pieces, retaining one piece for the next split and recording the other. This gives pairwise disjoint nonempty sets $(E_n)$.

The map
\[
S\subseteq\N\longmapsto\bigcup_{n\in S}E_n
\]
is injective from $\pow(\N)$ into the sigma-algebra.
\end{proof}

\begin{prop}
An algebra $\cA$ is a sigma-algebra if and only if it is closed under countable increasing unions.
\end{prop}

\begin{proof}
Only the converse needs proof. Given $A_n\in\cA$, put $B_n=\bigcup_{k=1}^nA_k$. Then $B_n\uparrow\bigcup_nA_n$, and the assumed closure gives $\bigcup_nA_n\in\cA$.
\end{proof}

\begin{prop}[Countable subgenerators]
For every family $\cE\subseteq\pow(X)$,
\[
\sigma(\cE)=\bigcup\{\sigma(\cF):\cF\subseteq\cE\text{ is countable}\}.
\]
\end{prop}

\begin{proof}
The right-hand side is closed under complements. If $A_n\in\sigma(\cF_n)$ and each $\cF_n$ is countable, then $\cF=\bigcup_n\cF_n$ is countable and all $A_n$ belong to $\sigma(\cF)$; hence so does $\bigcup_nA_n$. Thus the right-hand side is a sigma-algebra containing $\cE$, and equality follows.
\end{proof}

\section{Measures, null sets, and completion}

\begin{de}[Measure]
Let $(X,\cM)$ be a measurable space. A \emph{measure} is a function $\mu:\cM\to[0,\infty]$ such that $\mu(\varnothing)=0$ and
\[
\mu\left(\bigsqcup_{n=1}^{\infty}E_n\right)
 =\sum_{n=1}^{\infty}\mu(E_n)
\]
for every pairwise disjoint sequence $(E_n)$ in $\cM$. The triple $(X,\cM,\mu)$ is a \emph{measure space}.
\end{de}

\begin{eg}
Counting measure is given by $\#(E)=\card(E)$, with value $\infty$ for infinite sets. A Dirac measure at $x_0$ is $\delta_{x_0}(E)=\1_E(x_0)$. On an uncountable set with the countable/co-countable sigma-algebra, the formula
\[
\mu(E)=\begin{cases}
0,&E\text{ countable},\\
1,&E^c\text{ countable}
\end{cases}
\]
defines a probability measure.
\end{eg}

\begin{thm}[Basic properties]
Let $(X,\cM,\mu)$ be a measure space.
\begin{enumerate}
  \item If $E\subseteq F$, then $\mu(E)\leq\mu(F)$.
  \item For arbitrary $E_n\in\cM$,
  \[
  \mu\left(\bigcup_nE_n\right)\leq\sum_n\mu(E_n).
  \]
  \item If $E_n\uparrow E$, then $\mu(E_n)\uparrow\mu(E)$.
  \item If $E_n\downarrow E$ and $\mu(E_1)<\infty$, then $\mu(E_n)\downarrow\mu(E)$.
\end{enumerate}
\end{thm}

\begin{proof}
For monotonicity, write $F=E\sqcup(F\setminus E)$. For subadditivity, disjointify by setting
\[
F_1=E_1,\qquad F_n=E_n\setminus\bigcup_{k<n}E_k.
\]
Then $\bigcup_nF_n=\bigcup_nE_n$ and $\mu(F_n)\leq\mu(E_n)$.

For continuity from below, write $E$ as the disjoint union of $E_1$ and the increments $E_n\setminus E_{n-1}$. For continuity from above, set $F_n=E_1\setminus E_n$. Then $F_n\uparrow E_1\setminus E$, and subtraction is legitimate because $\mu(E_1)<\infty$:
\[
\mu(E_n)=\mu(E_1)-\mu(F_n)\longrightarrow
\mu(E_1)-\mu(E_1\setminus E)=\mu(E).
\]
\end{proof}

\begin{cor}[Liminf and limsup]
For every sequence $(E_n)$ in $\cM$,
\[
\mu(\liminf E_n)\leq\liminf_n\mu(E_n).
\]
If $\mu(\bigcup_nE_n)<\infty$, then
\[
\mu(\limsup E_n)\geq\limsup_n\mu(E_n).
\]
\end{cor}

\begin{proof}
The tail intersections $\bigcap_{k\geq n}E_k$ increase to $\liminf E_n$, and each has measure at most $\inf_{k\geq n}\mu(E_k)$. For the second inequality, the tail unions $\bigcup_{k\geq n}E_k$ decrease to $\limsup E_n$ and all lie in the finite-measure set $\bigcup_nE_n$, so continuity from above applies.
\end{proof}

\begin{prop}[Finite additivity criterion]
Let $\mu$ be finitely additive on an algebra $\cA$. Then $\mu$ is a premeasure if and only if it is continuous from below: whenever $E_n\uparrow E$ with all sets in $\cA$, one has $\mu(E_n)\uparrow\mu(E)$.
\end{prop}

\begin{proof}
A premeasure has continuity from below by the usual increment argument. Conversely, if $(E_n)$ is disjoint and $E=\bigcup_nE_n\in\cA$, then $F_N=\bigcup_{n\leq N}E_n$ increases to $E$, and finite additivity gives
\[
\mu(E)=\lim_N\mu(F_N)=\lim_N\sum_{n=1}^N\mu(E_n).
\]
\end{proof}

\begin{de}[Null sets and completeness]
A measurable set $N$ is \emph{null} if $\mu(N)=0$. The measure space is \emph{complete} if every subset of every measurable null set is measurable.
\end{de}

\begin{thm}[Completion]
Let $(X,\cM,\mu)$ be a measure space and define
\[
\overline{\cM}
 =\{A\cup N:A\in\cM,\ N\subseteq Z
      \text{ for some }Z\in\cM\text{ with }\mu(Z)=0\}.
\]
Then $\overline{\cM}$ is a sigma-algebra and
\[
\overline\mu(A\cup N)=\mu(A)
\]
is a well-defined complete measure on $\overline{\cM}$ extending $\mu$. It is the unique complete extension of $\mu$ whose measurable sets have this form.
\end{thm}

\begin{proof}
The family is closed under countable unions because a countable union of subsets of null sets is contained in a measurable null set. If $E=A\cup N$ with $N\subseteq Z$ and $\mu(Z)=0$, then
\[
E^c=(A^c\setminus Z)\cup\bigl((A^c\cap Z)\setminus N\bigr).
\]
The first term is measurable and the second is contained in the measurable null set $Z$, so the complement has the required form.

If $A\cup N=B\cup M$, with $N$ and $M$ contained in measurable null sets, then $A\triangle B$ is null, so $\mu(A)=\mu(B)$. Hence $\overline\mu$ is well defined. Disjoint representatives can be refined inside the original sigma-algebra; their overlaps are null, and countable additivity follows from countable additivity of $\mu$. Finally, every subset of a $\overline\mu$-null set is contained in an original measurable null set, proving completeness.
\end{proof}

\section{Sigma-finiteness, semifiniteness, and saturation}

\begin{de}
A measure is \emph{finite} if $\mu(X)<\infty$, \emph{sigma-finite} if $X$ is a countable union of finite-measure measurable sets, and \emph{semifinite} if every measurable $E$ with $\mu(E)=\infty$ contains a measurable $F$ with $0<\mu(F)<\infty$.
\end{de}

\begin{prop}
Every sigma-finite measure is semifinite.
\end{prop}

\begin{proof}
Choose increasing measurable sets $X_n\uparrow X$ with $\mu(X_n)<\infty$. If $\mu(E)=\infty$, then $\mu(E\cap X_n)\uparrow\infty$. Hence some $E\cap X_n$ has positive finite measure.
\end{proof}

\begin{prop}
If $\mu$ is semifinite, $\mu(E)=\infty$, and $C>0$, then there is a measurable $F\subseteq E$ with
\[
C<\mu(F)<\infty.
\]
\end{prop}

\begin{proof}
Let
\[
S=\sup\{\mu(F):F\subseteq E,\ F\in\cM,\ \mu(F)<\infty\}.
\]
If $S<\infty$, choose $F_n$ with $\mu(F_n)>S-1/n$ and set $G_n=\bigcup_{k\leq n}F_k$. Each $G_n$ has finite measure and therefore $\mu(G_n)\leq S$. Thus $G=\bigcup_nG_n$ has measure $S$. Since $\mu(E)=\infty$, the set $E\setminus G$ has infinite measure. Semifiniteness supplies a measurable $H\subseteq E\setminus G$ with $0<\mu(H)<\infty$, and then $G\cup H$ is a finite-measure subset of $E$ with measure larger than $S$, a contradiction. Hence $S=\infty$.
\end{proof}

\begin{prop}[Semifinite part]
For an arbitrary measure $\mu$, define
\[
\mu_{\mathrm{sf}}(E)
 =\sup\{\mu(F):F\subseteq E,\ F\in\cM,\ \mu(F)<\infty\}.
\]
Then $\mu_{\mathrm{sf}}$ is a semifinite measure and $\mu_{\mathrm{sf}}\leq\mu$. Moreover, $\mu_{\mathrm{sf}}=\mu$ if and only if $\mu$ is semifinite. If
\[
\nu(E)=
\begin{cases}
0,&E\text{ is a countable union of finite-$\mu$-measure sets},\\
\infty,&\text{otherwise},
\end{cases}
\]
then $\nu$ is a measure and $\mu=\mu_{\mathrm{sf}}+\nu$.
\end{prop}

\begin{proof}
Let $(E_n)$ be disjoint. For every finite $N$, finite-measure subsets $F_n\subseteq E_n$ give
\[
\mu_{\mathrm{sf}}\left(\bigcup_nE_n\right)
 \geq\sum_{n=1}^N\mu(F_n).
\]
Taking suprema and then $N\to\infty$ yields the lower inequality for countable additivity. Conversely, if $F\subseteq\bigcup_nE_n$ and $\mu(F)<\infty$, then
\[
\mu(F)=\sum_n\mu(F\cap E_n)
 \leq\sum_n\mu_{\mathrm{sf}}(E_n).
\]
Taking the supremum over $F$ gives the reverse inequality.

Semifiniteness of $\mu_{\mathrm{sf}}$ follows directly from the definition. If $\mu$ is semifinite, the preceding proposition shows that an infinite-measure set has finite-measure subsets of arbitrarily large measure, so $\mu_{\mathrm{sf}}(E)=\mu(E)$ for every $E$.

The class of countable unions of finite-$\mu$-measure sets is a sigma-ideal. Hence the displayed $\{0,\infty\}$-valued function $\nu$ is countably additive: a countable union fails to be sigma-finite exactly when at least one member of a disjoint decomposition fails to be sigma-finite. If $E$ is sigma-finite, then $\mu_{\mathrm{sf}}(E)=\mu(E)$ and $\nu(E)=0$; if it is not sigma-finite, then $\mu(E)=\nu(E)=\infty$. Thus $\mu=\mu_{\mathrm{sf}}+\nu$.
\end{proof}

\begin{de}[Local measurability and saturation]
A set $E\subseteq X$ is \emph{locally measurable} if
\[
E\cap A\in\cM
\quad\text{for every }A\in\cM\text{ with }\mu(A)<\infty.
\]
Let $\cM_{\mathrm{loc}}$ denote the collection of locally measurable sets. The measure space is \emph{saturated} if $\cM_{\mathrm{loc}}=\cM$.
\end{de}

\begin{prop}
The collection $\cM_{\mathrm{loc}}$ is a sigma-algebra containing $\cM$. If $\mu$ is sigma-finite, then $\mu$ is saturated.
\end{prop}

\begin{proof}
Closure of $\cM_{\mathrm{loc}}$ under complements and countable unions follows after intersecting with an arbitrary finite-measure measurable set. If $X_n\uparrow X$ and $\mu(X_n)<\infty$, then for $E\in\cM_{\mathrm{loc}}$,
\[
E=\bigcup_n(E\cap X_n)\in\cM.
\]
\end{proof}

\begin{prop}[Saturated extension]
Assume that $\mu$ is semifinite. For $E\in\cM_{\mathrm{loc}}$, set
\[
\widetilde\mu(E)
 =\sup\{\mu(A):A\in\cM,\ A\subseteq E,\ \mu(A)<\infty\}.
\]
Then $\widetilde\mu$ is a saturated measure on $\cM_{\mathrm{loc}}$ extending $\mu$.
\end{prop}

\begin{proof}
The proof of countable additivity is the same two-inequality argument used for the semifinite part: finite-measure measurable subsets of a disjoint union decompose into measurable intersections with the locally measurable summands. Semifiniteness shows that $\widetilde\mu(E)=\mu(E)$ for $E\in\cM$.

Suppose $E$ is locally measurable with respect to $(\cM_{\mathrm{loc}},\widetilde\mu)$. If $A\in\cM$ has finite $\mu$-measure, then $A\in\cM_{\mathrm{loc}}$ and $\widetilde\mu(A)<\infty$, so $E\cap A\in\cM_{\mathrm{loc}}$. Intersecting once more with $A$ shows $E\cap A\in\cM$. Therefore $E\in\cM_{\mathrm{loc}}$, proving saturation.
\end{proof}

\begin{remark}
Without semifiniteness, the extension that assigns $\infty$ to every locally measurable set outside $\cM$ can differ from the supremum construction. Here is a concrete example. Let $X=X_1\sqcup X_2$, where both sets are uncountable, and let
\[
\cM=\{A_1\cup A_2:A_1\subseteq X_1,\ A_2\subseteq X_2
       \text{ is countable or co-countable}\}.
\]
Define
\[
\mu(A_1\cup A_2)=\#A_1+\nu_2(A_2),
\qquad
\nu_2(A_2)=
\begin{cases}
0,&A_2\text{ countable},\\
\infty,&A_2\text{ co-countable}.
\end{cases}
\]
Every subset of $X_2$ is locally measurable, because every finite-measure measurable set meets $X_2$ in a countable set. Choose $E\subseteq X_2$ such that both $E$ and $X_2\setminus E$ are uncountable. Then $E\notin\cM$. The extension assigning $\infty$ outside $\cM$ gives $E$ value $\infty$, whereas the supremum over finite-measure measurable subsets of $E$ is $0$.
\end{remark}

\section{Outer measures and the Carath\'eodory construction}

\begin{de}[Outer measure]
An \emph{outer measure} on $X$ is a function $\mu^*:\pow(X)\to[0,\infty]$ satisfying
\begin{enumerate}
  \item $\mu^*(\varnothing)=0$;
  \item $A\subseteq B$ implies $\mu^*(A)\leq\mu^*(B)$;
  \item $\mu^*(\bigcup_nA_n)\leq\sum_n\mu^*(A_n)$.
\end{enumerate}
\end{de}

\begin{prop}[Outer measure from coverings]
Let $\cE\subseteq\pow(X)$ contain $\varnothing$ and $X$, and let $p:\cE\to[0,\infty]$ satisfy $p(\varnothing)=0$. Define
\[
\mu^*(A)=\inf\left\{\sum_{j=1}^{\infty}p(E_j):
 E_j\in\cE,\ A\subseteq\bigcup_jE_j\right\}.
\]
Then $\mu^*$ is an outer measure.
\end{prop}

\begin{proof}
Only countable subadditivity requires work. If $\sum_i\mu^*(A_i)=\infty$, there is nothing to prove. Otherwise, given $\varepsilon>0$, choose a cover $(E_{ij})_j$ of $A_i$ with cost below $\mu^*(A_i)+\varepsilon2^{-i}$. The family of all $E_{ij}$ covers $\bigcup_iA_i$, so
\[
\mu^*\left(\bigcup_iA_i\right)
 \leq\sum_i\mu^*(A_i)+\varepsilon.
\]
Let $\varepsilon\downarrow0$.
\end{proof}

\begin{de}[Carath\'eodory measurability]
A set $A\subseteq X$ is \emph{$\mu^*$-measurable} if
\[
\mu^*(E)=\mu^*(E\cap A)+\mu^*(E\cap A^c)
\quad\text{for every }E\subseteq X.
\]
\end{de}

\begin{thm}[Carath\'eodory]
The collection $\cM^*$ of $\mu^*$-measurable sets is a sigma-algebra. The restriction $\mu^*\restr\cM^*$ is a complete measure.
\end{thm}

\begin{proof}
Complements are immediate. To prove closure under countable unions, first disjointify a sequence $(A_n)$ by
\[
C_1=A_1,\qquad C_n=A_n\setminus\bigcup_{k<n}A_k.
\]
Finite iteration of the Carath\'eodory identity gives
\[
\mu^*(E)\geq
 \sum_{n=1}^N\mu^*(E\cap C_n)
 +\mu^*\left(E\setminus\bigcup_{n=1}^NC_n\right).
\]
Letting $N\to\infty$ and using subadditivity yields the equality required for $\bigcup_nA_n$. Setting $E=\bigcup_nC_n$ proves countable additivity on disjoint measurable sets.

If $N$ has outer measure zero, then every subset $A\subseteq N$ has outer measure zero and satisfies the Carath\'eodory identity. Thus the measure is complete.
\end{proof}

\begin{de}[Premeasure]
Let $\cA$ be an algebra. A function $\mu_0:\cA\to[0,\infty]$ is a \emph{premeasure} if $\mu_0(\varnothing)=0$ and
\[
\mu_0\left(\bigsqcup_nA_n\right)=\sum_n\mu_0(A_n)
\]
whenever the $A_n$ are disjoint members of $\cA$ and their union belongs to $\cA$.
\end{de}

\begin{thm}[Extension from a premeasure]
Let $\mu_0$ be a premeasure on an algebra $\cA$. Define
\[
\mu^*(E)=\inf\left\{\sum_n\mu_0(A_n):
 A_n\in\cA,\ E\subseteq\bigcup_nA_n\right\}.
\]
Then
\begin{enumerate}
  \item $\mu^*(A)=\mu_0(A)$ for every $A\in\cA$;
  \item every $A\in\cA$ is $\mu^*$-measurable;
  \item $\mu=\mu^*\restr\sigma(\cA)$ is a measure extending $\mu_0$.
\end{enumerate}
If $\mu_0$ is sigma-finite, this extension is unique on $\sigma(\cA)$.
\end{thm}

\begin{proof}
The cover consisting only of $A$ gives $\mu^*(A)\leq\mu_0(A)$. Conversely, if $A\subseteq\bigcup_nA_n$, disjointify the cover inside $A$:
\[
B_n=(A\cap A_n)\setminus\bigcup_{k<n}A_k.
\]
Then $A=\bigsqcup_nB_n$, so
\[
\mu_0(A)=\sum_n\mu_0(B_n)\leq\sum_n\mu_0(A_n).
\]
Taking the infimum gives equality.

For $A\in\cA$, split every covering set $B_n$ into $B_n\cap A$ and $B_n\cap A^c$. Finite additivity of $\mu_0$ and the definition of $\mu^*$ yield
\[
\mu^*(E)\geq\mu^*(E\cap A)+\mu^*(E\cap A^c).
\]
The reverse inequality is subadditivity, so $A$ is measurable. Existence follows from Carath\'eodory's theorem.

For uniqueness, first assume $\mu_0(X)<\infty$. If $\nu$ is another extension, the class
\[
\cL=\{E\in\sigma(\cA):\mu(E)=\nu(E)\}
\]
is a Dynkin system containing the pi-system $\cA$, hence contains $\sigma(\cA)$. In the sigma-finite case, partition $X$ into disjoint sets in $\cA$ of finite premeasure and apply the finite argument to the restrictions before summing.
\end{proof}

\begin{prop}[Measurable hulls]
Let $\cA_\sigma$ be the collection of countable unions of members of $\cA$.
\begin{enumerate}
  \item For every $E\subseteq X$ and $\varepsilon>0$, there is $A\in\cA_\sigma$ with $E\subseteq A$ and
  \[
  \mu^*(A)\leq\mu^*(E)+\varepsilon.
  \]
  \item If $\mu^*(E)<\infty$, there is a measurable $B\supseteq E$ with $\mu^*(B)=\mu^*(E)$. If, in addition, $E\in\cM^*$, then $\mu^*(B\setminus E)=0$.
  \item Under sigma-finiteness, every $E\in\cM^*$ has a measurable hull $B\supseteq E$ satisfying $\mu^*(B\setminus E)=0$, even when $\mu^*(E)=\infty$.
\end{enumerate}
\end{prop}

\begin{proof}
The first assertion is the definition of the covering outer measure. If $\mu^*(E)<\infty$, choose $A_n\in\cA_\sigma$ containing $E$ with $\mu^*(A_n)\leq\mu^*(E)+2^{-n}$, and replace $A_n$ by $\bigcap_{k\leq n}A_k$ to make the sequence decreasing. Then $B=\bigcap_nA_n$ is measurable, contains $E$, and continuity from above inside the finite-measure set $A_1$ gives $\mu^*(B)=\mu^*(E)$. If $E$ is measurable, then
\[
\mu^*(B)=\mu^*(E)+\mu^*(B\setminus E),
\]
so the difference is null.

Under sigma-finiteness, take a disjoint measurable partition $X=\bigsqcup_nX_n$ with $\mu^*(X_n)<\infty$. Apply the finite construction to $E\cap X_n$, obtaining measurable $B_n\subseteq X_n$ with $E\cap X_n\subseteq B_n$ and $\mu^*(B_n\setminus E)=0$. Then $B=\bigcup_nB_n$ has the required properties.
\end{proof}

\begin{de}[Inner measure on a finite space]
If $\mu^*(X)<\infty$, define
\[
\mu_*(E)=\mu^*(X)-\mu^*(X\setminus E).
\]
\end{de}

\begin{prop}
When $\mu^*(X)<\infty$, a set $E\subseteq X$ is $\mu^*$-measurable if and only if $\mu^*(E)=\mu_*(E)$.
\end{prop}

\begin{proof}
Measurability immediately gives equality by applying the Carath\'eodory identity to $X$. Conversely, suppose equality holds and let $B\supseteq E$ be a measurable hull with $\mu^*(B)=\mu^*(E)$. Since $X\setminus B$ is measurable,
\[
\mu^*(X\setminus E)
 =\mu^*(X\setminus B)+\mu^*(B\setminus E).
\]
The assumed equality forces $\mu^*(B\setminus E)=0$. Completeness then gives $E=B\setminus(B\setminus E)\in\cM^*$.
\end{proof}

\begin{prop}[Regenerating an outer measure]
Let $\mu=\mu^*\restr\cM^*$ and define
\[
\mu^+(E)=\inf\{\mu(A):A\in\cM^*,\ E\subseteq A\}.
\]
Then $\mu^*(E)\leq\mu^+(E)$ for every $E$. Equality holds precisely when $E$ has a measurable hull of the same outer measure.
\end{prop}

\begin{proof}
Every measurable cover $A\supseteq E$ satisfies $\mu^*(E)\leq\mu(A)$, so $\mu^*(E)\leq\mu^+(E)$. If $E$ has a measurable hull $B$ with $\mu(B)=\mu^*(E)$, then $\mu^+(E)\leq\mu(B)$ and equality follows. Conversely, if $\mu^+(E)=\mu^*(E)<\infty$, choose measurable $A_n\supseteq E$ with $\mu(A_n)\leq\mu^*(E)+2^{-n}$, make the sequence decreasing, and put $B=\bigcap_nA_n$; continuity from above gives $\mu(B)=\mu^*(E)$. If the common value is infinite, every measurable superset of $E$ already has measure $\infty$, so any such superset, for example $X$, is a hull of the same outer measure.
\end{proof}

\begin{eg}
Let $X=\{a,b\}$ and define
\[
\mu^*(\varnothing)=0,\quad
\mu^*(\{a\})=2,\quad
\mu^*(\{b\})=3,\quad
\mu^*(X)=4.
\]
The only Carath\'eodory-measurable sets are $\varnothing$ and $X$. Hence the regenerated outer measure is $4$ on every nonempty set and differs from $\mu^*$ on the singletons.
\end{eg}

\begin{prop}[Saturation of the induced measure]
The complete measure $\mu^*\restr\cM^*$ obtained from a premeasure is saturated.
\end{prop}

\begin{proof}
Let $E$ be locally measurable. If $\mu^*(E)<\infty$, choose a finite measurable hull $B\supseteq E$. Then $E=E\cap B$ is measurable by local measurability. If $\mu^*(E)=\infty$ and $E$ were not measurable, the Carath\'eodory identity would fail for some test set $A$. The strict failure forces $\mu^*(A)<\infty$. Replacing $A$ by a measurable hull of the same finite outer measure preserves the strict failure, but local measurability makes both intersections with that hull measurable, a contradiction.
\end{proof}

\begin{prop}[Completion and the Carath\'eodory domain]
If the premeasure $\mu_0$ is sigma-finite, the Carath\'eodory measure space is the completion of the measure on $\sigma(\cA)$. Without sigma-finiteness, every finite-measure Carath\'eodory-measurable set belongs to that completion, and the full Carath\'eodory space is the saturation of the completion.
\end{prop}

\begin{proof}
Under sigma-finiteness, every Carath\'eodory-measurable set has a measurable hull in $\sigma(\cA)$ whose difference is null, so it belongs to the completion. The same argument works for every finite-measure set without sigma-finiteness. Every Carath\'eodory-measurable set is locally measurable with respect to the completion, and saturation then recovers the remaining sets.
\end{proof}

\begin{prop}[Measure on traces]
Let $(X,\cM,\mu)$ be finite, let $\mu^*$ be the outer measure induced by $\mu$, and let $E\subseteq X$ satisfy $\mu^*(E)=\mu(X)$. Put
\[
\cM_E=\{A\cap E:A\in\cM\},
\qquad
\nu(A\cap E)=\mu(A).
\]
Then $\nu$ is a well-defined measure on $E$.
\end{prop}

\begin{proof}
For every $A\in\cM$,
\[
\mu(X)=\mu^*(E)
 \leq\mu^*(E\cap A)+\mu^*(E\cap A^c)
 \leq\mu(A)+\mu(A^c)=\mu(X).
\]
Therefore $\mu^*(E\cap A)=\mu(A)$. If $A\cap E=B\cap E$, then $E\cap(A\triangle B)=\varnothing$, so $\mu(A\triangle B)=0$ and $\mu(A)=\mu(B)$. This proves well-definedness. For disjoint traces, measurable representatives can be disjointified without changing their intersections with $E$, after which countable additivity follows from that of $\mu$.
\end{proof}

\section{Lebesgue--Stieltjes and Lebesgue measure}

\begin{de}[Half-open intervals]
An $h$-interval is a set of the form $(a,b]$, where $-\infty\leq a<b\leq\infty$, with the evident conventions at infinite endpoints. Finite disjoint unions of $h$-intervals form an algebra $\cA$ on $\R$.
\end{de}

\begin{thm}[Lebesgue--Stieltjes premeasure]
Let $F:\R\to\R$ be increasing and right-continuous. Set
\[
\mu_0((a,b])=F(b)-F(a)
\]
for finite endpoints and use the monotone limits of $F$ at $\pm\infty$ for unbounded intervals. Extend $\mu_0$ additively to finite disjoint unions of $h$-intervals. Then $\mu_0$ is a premeasure on $\cA$.
\end{thm}

\begin{proof}
A common refinement by all endpoints shows that the value assigned to a finite union is independent of its representation; the sums telescope on each refined interval.

Suppose first that a bounded interval $(a,b]$ is the disjoint union of intervals $(a_j,b_j]$. Finite partial unions give
\[
\mu_0((a,b])\geq\sum_{j=1}^N\mu_0((a_j,b_j]).
\]
For the reverse inequality, fix $\varepsilon>0$. Right continuity permits a slight rightward enlargement of each covering interval with total added $F$-increment below $\varepsilon$, and a slight truncation near $a$ with error below $\varepsilon$. The enlarged open intervals cover a compact interval, so a finite subcover exists. Ordering this finite cover and using telescoping yields
\[
\mu_0((a,b])\leq\sum_j\mu_0((a_j,b_j])+2\varepsilon.
\]
Let $\varepsilon\downarrow0$. Finite unions and infinite endpoints follow by refinement and monotone limits.
\end{proof}

\begin{thm}[Lebesgue--Stieltjes measure]
There is a unique Borel measure $\mu_F$ on $\R$, finite on bounded sets, such that
\[
\mu_F((a,b])=F(b)-F(a).
\]
Conversely, every Borel measure on $\R$ that is finite on bounded sets is of this form for an increasing right-continuous function $F$, unique up to an additive constant.
\end{thm}

\begin{proof}
Existence and uniqueness follow from the extension theorem and sigma-finiteness of the interval algebra. Conversely, given $\mu$, define
\[
F(x)=\begin{cases}
\mu((0,x]),&x\geq0,\\
-\mu((x,0]),&x<0.
\end{cases}
\]
Continuity from above and below gives right continuity, and the interval identity determines $\mu$.
\end{proof}

\begin{thm}[Regularity on $\R$]
Let $\mu_F$ be a Lebesgue--Stieltjes measure and let $E$ be measurable in its completion. Then
\[
\mu_F(E)=\inf\{\mu_F(U):E\subseteq U,\ U\text{ open}\}
\]
after the usual interpretation when the value is infinite. If $E$ has finite measure, then
\[
\mu_F(E)=\sup\{\mu_F(K):K\subseteq E,\ K\text{ compact}\}.
\]
Every measurable set differs from a $G_\delta$ set and from an $F_\sigma$ set by a null set.
\end{thm}

\begin{proof}
Open-interval covers and the defining outer measure give outer regularity. For a bounded finite-measure $E$, choose a compact interval $I$ containing $E$ and an open set $U\supseteq I\setminus E$ with
\[
\mu_F(U)\leq\mu_F(I\setminus E)+\varepsilon.
\]
Then $K=I\setminus U$ is compact, $K\subseteq E$, and
\[
\mu_F(K)\geq\mu_F(E)-\varepsilon.
\]
An increasing bounded exhaustion handles unbounded sets.

For the descriptive-set statements, choose open supersets with summable excess, make them nested, and intersect them to obtain a $G_\delta$ hull with null difference. Apply the same construction to the complement to obtain the $F_\sigma$ representation. A disjoint finite-measure exhaustion handles the sigma-finite case without relying on an invalid assertion that an arbitrary countable union of $G_\delta$ sets is again $G_\delta$.
\end{proof}

\begin{de}[Lebesgue measure]
Taking $F(x)=x$ gives Lebesgue measure $m$ on the Borel sets of $\R$. Its completion is also denoted by $m$, and its domain is the sigma-algebra of Lebesgue-measurable sets.
\end{de}

\begin{thm}[Translation and dilation]
If $E$ is Lebesgue measurable, $t\in\R$, and $r\neq0$, then $E+t$ and $rE$ are Lebesgue measurable and
\[
m(E+t)=m(E),
\qquad
m(rE)=|r|m(E).
\]
\end{thm}

\begin{proof}
First work on Borel sets. The measures $E\mapsto m(E+t)$ and $E\mapsto |r|^{-1}m(rE)$ agree with $m$ on half-open intervals, so uniqueness of the sigma-finite extension gives equality on all Borel sets. Borel null sets are therefore sent to null sets, and the identities pass to the completion.
\end{proof}

\begin{prop}[The Cantor set]
The middle-thirds Cantor set $C\subseteq[0,1]$ is compact, perfect, totally disconnected, and has Lebesgue measure zero. The Cantor function is continuous and nondecreasing, is constant on each component of $C^c$, and maps $[0,1]$ onto $[0,1]$.
\end{prop}

\begin{proof}
At stage $n$, the construction leaves $2^n$ intervals of length $3^{-n}$, whose total length is $(2/3)^n$. Compactness follows from nested closedness, and the remaining topological properties follow from the ternary construction.
\end{proof}

\section{Consequences and corrected exercises}

\begin{thm}[Positive-measure sets contain nonmeasurable subsets]
Every Lebesgue-measurable set $E\subseteq\R$ with $m(E)>0$ contains a nonmeasurable subset.
\end{thm}

\begin{proof}
Replace $E$ by a bounded positive-measure subset. On $E$, declare $x\sim y$ when $x-y\in\Q$, and choose one representative from each equivalence class, obtaining $V\subseteq E$. Distinct rational translates of $V$ are disjoint. If $V$ were measurable with positive measure, infinitely many translates $V+q$ with $q$ in a bounded rational interval would be pairwise disjoint and contained in one bounded interval, which is impossible. Hence $m(V)=0$. But $E$ is contained in the union of countably many rational translates of $V$, which would then be null, a contradiction. Thus $V$ is nonmeasurable.
\end{proof}

\begin{prop}[High density in an interval]
If $E$ is Lebesgue measurable and $m(E)>0$, then for every $\alpha<1$ there is an open interval $I$ such that
\[
m(E\cap I)>\alpha m(I).
\]
\end{prop}

\begin{proof}
For $\alpha\leq0$, choose an interval $I$ with $m(E\cap I)>0$; the conclusion is immediate. Thus assume $0<\alpha<1$ and replace $E$ by a bounded positive-measure subset. If the conclusion failed, every interval would satisfy $m(E\cap I)\leq\alpha m(I)$. Choose an open-interval cover $(I_j)$ of $E$ with
\[
\sum_jm(I_j)<\frac{m(E)}{\alpha}.
\]
Then
\[
m(E)\leq\sum_jm(E\cap I_j)
 \leq\alpha\sum_jm(I_j)<m(E),
\]
a contradiction.
\end{proof}

\begin{thm}[Steinhaus]
If $E\subseteq\R$ is Lebesgue measurable and $m(E)>0$, then the difference set
\[
E-E=\{x-y:x,y\in E\}
\]
contains an open interval centered at $0$.
\end{thm}

\begin{proof}
Choose a compact set $K\subseteq E$ with $m(K)>0$, and then an open set $U\supseteq K$ with $m(U)<2m(K)$. Compactness gives $\delta>0$ such that $K+h\subseteq U$ whenever $|h|<\delta$. If $K$ and $K+h$ were disjoint, translation invariance would give
\[
m(U)\geq m(K\cup(K+h))=2m(K),
\]
a contradiction. Hence $K\cap(K+h)\neq\varnothing$, so $h\in K-K\subseteq E-E$.
\end{proof}

\begin{prop}
There exists a Borel set $A\subseteq[0,1]$ such that for every nonempty interval $I\subseteq[0,1]$,
\[
0<m(A\cap I)<m(I).
\]
\end{prop}

\begin{proof}
Enumerate the open intervals with rational endpoints contained in $[0,1]$ as $(I_n)$. Recursively choose disjoint compact nowhere-dense sets $K_n,L_n\subseteq I_n$, each of positive measure, and disjoint from all previously chosen sets. This is possible because a finite union of compact nowhere-dense sets has open dense complement in $I_n$, inside which one can place two small positive-measure fat Cantor sets.

Set $A=\bigcup_nK_n$. Then $A$ is $F_\sigma$, hence Borel. Every nonempty interval contains some $I_n$; the set $K_n$ gives positive measure inside $A$, while $L_n$ remains disjoint from every $K_j$ and gives positive measure inside $A^c$.
\end{proof}
